\begin{textsample}{POS Dim 9 – global\_south\_2025\_06 – Score 1.00 – global\_south\_2025\_06/125544251.txt}  \label{ex:f9_pos_176}
U.S cast lone vote against permanent ceasefire in Gaza The United States ( U.S ) vetoed a UN Security Council resolution on Wednesday, June 5, 2025, that called for an immediate, unconditional, and permanent ceasefire in Gaza, as well as unrestricted humanitarian aid access.
The resolution, which received 14 votes in favor, aimed to address the catastrophic humanitarian situation in Gaza, where famine looms and aid has been limited since Israel lifted an 11-week blockade.
The text was co-sponsored by Algeria, Denmark, Greece, Guyana, Pakistan, Panama, the Republic of Korea, Sierra Leone, Slovenia and Somalia, collectively known as the E-10.
The text received 14 votes in favour, with the U.S. casting the lone vote against, blocking the initiative backed by all 10 elected members of the Council.
As one of the Council ' s five permanent members, the U.S. holds veto power, a negative vote that automatically blocks any resolution from going forward.
If the resolution had @ @ @ @ @ @ @ @ @ @ immediate, unconditional and permanent ceasefire in Gaza"to be respected by all parties.
The text reaffirmed the Council ' s earlier call for the"immediate, dignified and unconditional release of all hostages held by Hamas and other groups."The draft expressed grave concern over the"catastrophic humanitarian situation"in Gaza, following months of almost total Israeli aid blockade, including the risk of famine, highlighted by recent international food security assessments.
The draft resolution reaffirmed the obligation of all parties to comply with international law, including international humanitarian and human rights law.
In addition to a ceasefire, the draft resolution demanded the"immediate and unconditional lifting of all restrictions"on the entry and distribution of humanitarian aid in Gaza, It called for safe and \textbf{unhindered} access for UN and humanitarian partners across the enclave.
It also urged the restoration of essential services, in accordance with humanitarian principles and prior Security Council resolutions.
The text voiced support for ongoing mediation efforts led by Egypt @ @ @ @ @ @ @ @ @ @ framework outlined in resolution 2735 ( 2024 ).
Resolution 2735 had envisioned a permanent cessation of hostilities, the release of all hostages, the exchange of Palestinian prisoners, the return of all remains, full Israeli military withdrawal from Gaza and the start of a long-term reconstruction plan.
Speaking ahead of the vote, acting U.S.
Representative Dorothy Shea described the draft resolution as"unacceptable"."U.S. opposition to this resolution should come as no surprise, it is unacceptable for what it does say, it is unacceptable for what it does not say, and it is unacceptable for the manner in which it has been advanced,"she said."The United States has been clear, we would not support any measure that fails to condemn Hamas and does not call for Hamas to disarm and leave Gaza,"Shea stressed.
She added that Hamas had rejected numerous ceasefire proposals, including one over the weekend that would have provided a pathway to end the conflict and release the remaining @ @ @ @ @ @ @ @ @ @ Council to award Hamas ' intransigence,"Shea said, stressing,"Hamas and other terrorists must have no future in Gaza."As Secretary ( Marco ) Rubio has said : ' If an ember survives, it will spark again into a fire '."The failure of the resolution comes as the humanitarian crisis in Gaza deepens, with UN agencies warning of the total collapse of health services, growing displacement, and a rising death toll around the new privatised U.S.-Israel led aid distribution system which bypasses established agencies."The world is watching, day after day, horrifying scenes of Palestinians being shot, wounded or killed in Gaza? while simply trying to eat,"UN relief chief Tom Fletcher, said ( NAN )

% matched lemmas: unhindered
\end{textsample}
